
\documentclass[20pt, a4paper]{extarticle}

% ─── Paquetes ────────────────────────────────────────────────────────────────
\usepackage[utf8]{inputenc}
\usepackage[T1]{fontenc}
\usepackage[spanish]{babel}
\usepackage{geometry}
\usepackage{graphicx}
\usepackage{float}
\usepackage{hyperref}
\usepackage{xcolor}
\usepackage{amsmath}
\usepackage{fancyhdr}
\usepackage{titlesec}
\usepackage{enumitem}
\usepackage{setspace}
\usepackage{parskip}
\usepackage{lmodern}
\usepackage{microtype}
\usepackage{tikz}
\usepackage{mdframed}

% ─── Geometría ───────────────────────────────────────────────────────────────
\geometry{
  top=3cm, bottom=3cm,
  left=1cm, right=1cm,
  headheight=15pt
}

% ─── Colores ─────────────────────────────────────────────────────────────────
\definecolor{azulUR}{RGB}{0, 60, 120}
\definecolor{dorado}{RGB}{180, 140, 50}
\definecolor{grisclaro}{RGB}{245, 245, 245}

% ─── Interlineado ────────────────────────────────────────────────────────────
\setstretch{1.45}
\setlength{\parskip}{0.8em}

% ─── Hyperref ────────────────────────────────────────────────────────────────
\hypersetup{
  colorlinks=true,
  linkcolor=azulUR,
  urlcolor=azulUR,
  hidelinks
}

% ─── Encabezado y pie de página ──────────────────────────────────────────────
\pagestyle{fancy}
\fancyhf{}
\renewcommand{\headrulewidth}{0.4pt}
\renewcommand{\footrulewidth}{0.4pt}
\fancyhead[C]{\small\color{azulUR}\textit{Discurso de Graduación · CCT · UR}}
\fancyfoot[C]{\small\color{gray}\thepage}

% ─── Estilo sección decorativa ───────────────────────────────────────────────
\newcommand{\separador}{%
  \begin{center}
    \textcolor{dorado}{\rule{2cm}{0.4pt}\quad\(\star\)\quad\rule{2cm}{0.4pt}}
  \end{center}
  \vspace{0.3em}
}

% ─── Entorno para citas destacadas ───────────────────────────────────────────
\newmdenv[
  backgroundcolor=grisclaro,
  linecolor=dorado,
  linewidth=2pt,
  topline=false,
  rightline=false,
  bottomline=false,
  innerleftmargin=14pt,
  innerrightmargin=10pt,
  innertopmargin=8pt,
  innerbottommargin=8pt,
  skipabove=1em,
  skipbelow=1em
]{citadestacada}

% ═════════════════════════════════════════════════════════════════════════════
\begin{document}


% ─── Cuerpo del discurso ─────────────────────────────────────────────────────
\thispagestyle{fancy}


Buenas, como ya me ha presentado Alexia soy Pablo Galilea y queremos cerrar estos agradecimientos recordando la importancia de la universidad pública, con una gratitud que no es solo emoción, sino conciencia. Porque sin sus aulas abiertas, sin ese sueño colectivo sostenido entre todas y todos, muchos jamás habríamos llegado hasta aquí. La universidad pública es el motor de la formación del pueblo y convierte el derecho a pensar en un territorio común. Por eso no basta con nombrarla: necesitamos defenderla, exigirla fuerte, digna y bien financiada. Y es urgente que comprendamos, de una vez y para siempre, que su valor no se mide en presupuestos, sino en las vidas que transforma y en la esperanza que siembra cada vez que abre la puerta a quien no tiene otra llave.

\separador

\textbf{\large\color{azulUR} El tiempo}

\medskip

Al principio del discurso os hemos dicho que os querimaos hablar de tres cosas, la primera eran las personas, como bien ha hecho Alexia, la segunda es el tiempo.

En este momento, somos jóvenes. Estamos sanos. Y somos potencialmente felices. Los tres a la vez.
¿Sabéis lo raro que es eso? ¿Lo privilegiado qué es?

Durante la carrera hemos ido siempre con prisas.
Queremos terminar el semestre para llegar al siguiente, terminar el año para llegar al siguiente, terminar la carrera para empezar a trabajar, empezar a trabajar\ldots ¿para qué? ¿Para seguir corriendo?
\separador




Hay algo muy triste en vivir la vida como un pasillo que hay que cruzar lo más rápido posible.


Disfruta de tu vida. Haz lo que te gusta. Por supuesto que hay que formarse, hay que trabajar, hay que ser responsable. Pero hay tiempo para todo. Lo que no puede ser es que tu vida sea una carrera en la que solo enciendas el cerebro en pequeños descansos. Haz que esos descansos sean la mayor parte del día.

Pasa tiempo con tu familia. Pasa tiempo con tus amigos. Pasa tiempo con tu pareja. Permite que te pasen cosas.

Porque cuando hayas visto la vida pasar, no te vas a acordar del número de integrales que has resuelto.
Ni del número de líneas de código que has picado. Ni de las pérdidas de carga de las tuberías de hidráulica que has hecho. Ni de todas las fórmulas químicas que has memorizado. Ni de los vinos probados en la sala de catas.
Te acordarás de los momentos que has vivido con las personas que quieres y haciendo lo que te gusta. 

El éxito no es una nota. No es un sueldo.
El éxito es saber disfrutar de los pequeños momentos. Es tener la capacidad de encontrar la belleza en lo cotidiano. Es ser capaz de reírse de uno mismo. Es no tomarse la vida tan en serio, porque al final, lo que más importa es que te guste cómo vives, no cuánto ganas o qué título tienes.
En resumen, es irse a dormir con la sensación de haber \textit{vivido} el día, no de haberlo sobrevivido.

Sé consciente de lo privilegiado que eres.
No es un reproche, es un recordatorio.
Porque esta vida hay que disfrutarla, y nadie lo va a hacer por ti.

\separador

\textbf{\large\color{azulUR} La responsabilidad}

\medskip

Para concluir, llegamos a la piedra angular del discurso, el último de los 3 temas: la responsabilidad.

Y es que hay que ser buenos, además de disfrutar.

Somos las mentes del futuro. No solo de esta universidad, sino del mundo.
Y este, ahora mismo, necesita urgentemente mentes honestas.
Necesita gente que se indigna cuando algo es injusto.
Gente que no mira para otro lado.
Gente que pone su granito de arena aunque sienta que es pequeño, porque si todos lo hacemos, llegaremos a formar cosas enormes.

No podemos seguir normalizando que quienes tienen el poder para cambiar el mundo lo usen para destruirlo.
La capacidad de poder cambiarlo está en las decisiones que tomamos cada día: en cómo votamos, cómo consumimos, cómo hablamos y cómo tratamos a los demás.

No podemos permitir que la tecnología que hemos aprendido a dominar se use para espiar, para manipular, para controlar, para dividir. 

Todo mi apoyo y mi cariño para las víctimas de Palestina, de Irán, y de todos los conflictos que no aparecen en los titulares.
Porque el hecho de que no se hable de ellos no significa que no existan.



Si con estas palabras conseguimos que uno solo de vosotros salga de aquí pensando un poco más en cómo quiere vivir y qué mundo quiere construir, nos damos por satisfechos.

\vspace{1em}

{\large\bfseries Enhorabuena a todos. Ha sido un honor compartir este camino con vosotros.}
  \vspace{2.5cm}
\end{document}
